%%% ------------------------------------------------------------------
%%% JPC3 --- Improving the Czech Wordnet / Open Multilingual Wordnet
%%% Project report skeleton.
%%%
%%% This is a skeleton.  Replace everything marked TODO, delete the
%%% parts you do not need, and keep the parts you do.  Every project
%%% has the same four-part spine, whatever its topic:
%%%     Extend  ->  Validation  ->  Visualise  ->  Release
%%% so keep those four sections even if the rest changes.  The weight is
%%% on Validation: the extraction is scaffolded, the judging is the work.
%%%
%%% To compile on your own machine:
%%%     pdflatex report
%%%     bibtex   report
%%%     pdflatex report
%%%     pdflatex report
%%% (twice at the end, so the citations and cross-references settle)
%%%
%%% On Overleaf just press Recompile; it does all of that for you.
%%% ------------------------------------------------------------------

\documentclass[a4paper,12pt]{article}

%%% uniform margins
\usepackage[margin=25mm]{geometry}

\usepackage[utf8]{inputenc}       % delete this line if you compile with XeLaTeX
\usepackage[T1]{fontenc}          % proper accented letters: ě š č ř ž ý á í é
\usepackage[english]{babel}

%%% citations: \citep{...} and \citet{...}
\usepackage[round]{natbib}
\bibliographystyle{plainnat}

%%% tables and figures
\usepackage{booktabs}             % \toprule \midrule \bottomrule
\usepackage{graphicx}

%%% hyperlinks
\usepackage[svgnames]{xcolor}
\usepackage{hyperref,url}
\hypersetup{
    colorlinks=true,
    linkcolor=DarkBlue,
    urlcolor=DarkRed,
}

%%% numbered examples and glosses (you will need these for Czech data)
\usepackage{langsci-gb4e}
\let\eachwordone=\itshape         % italicises the source line of a gloss

%%% Some versions of gb4e need \noautomath so that _ and ^ behave
%%% outside maths; older ones do not have it at all, so provide a
%%% harmless fallback rather than let the file fail to compile.
\providecommand{\noautomath}{}
\noautomath

%%% A shorthand for glossing labels, so they come out as small capitals:
%%% type \gm{acc} and you get ACC.
\newcommand{\gm}[1]{\textsc{\MakeLowercase{#1}}}

\title{TODO: a title that names what you did\\
  \large JPC3 --- Improving the Czech Wordnet}
\author{TODO: First Author's Name \and TODO: Second Author's Name}
\date{2027-01-31}

\begin{document}
\maketitle

\begin{abstract}
  TODO: four or five sentences.  What part of the wordnet did you set
  out to improve, how did you do it (semi-automatically?), how much did
  you add, and how good was it?  Give the headline number here.  Write
  this last.
\end{abstract}

%%% ------------------------------------------------------------------
\section{Introduction}
%%% ------------------------------------------------------------------

TODO: Say which gap in the Czech wordnet you are filling and why it
matters.  Lead with a concrete example of something that is missing or
wrong, not with the literature.  State your \emph{numerical goal} up
front --- the target you agreed at the start --- so the reader knows
what ``done'' means.

Here is how you give a numbered example and refer back to it
as~(\ref{ex:gap}):

\begin{exe}
  \ex \label{ex:gap} TODO: a synset, lemma, or relation that your
  project adds or fixes.
\end{exe}

%%% ------------------------------------------------------------------
\section{Data and resources}
%%% ------------------------------------------------------------------

TODO: Name every resource you used --- the version of the Czech
wordnet, any corpus, lexicon, image set, or model --- and give its
licence.  A reader must be able to get the same inputs you had.

If you quote Czech examples they must be glossed, so that a reader who
does not speak Czech can follow.  The source line and the gloss line
must have the \emph{same number of words}, and matching hyphens, or the
columns will not line up:

\begin{exe}
  \ex Czech [ces] \\
      \gll učitel \quad\textrightarrow\quad učitel-ka \\
           teacher.\gm{m} \qquad teacher-\gm{f} \\
  \trans ``(male) teacher $\rightarrow$ (female) teacher''
\end{exe}

Name the language on the first line with its ISO 639-3 code in square
brackets: Czech \texttt{[ces]}, English \texttt{[eng]}.  Cite the
source of any example you did not make up yourself.

%%% ------------------------------------------------------------------
\section{Previous work}
%%% ------------------------------------------------------------------

TODO: What have other people done?  You need at least three references.
Read them --- do not cite a paper you have not opened.  The two starter
readings for your task are in \texttt{report.bib}; find the third
yourself.

Use \verb|\citet| when the author is part of your sentence, and
\verb|\citep| when the citation sits in brackets:
\citet{palasmrz2004} built the original Czech wordnet; it is one of
many in the Open Multilingual Wordnet \citep{bond-foster2013}.

%%% ------------------------------------------------------------------
\section{Method and results}
%%% ------------------------------------------------------------------

TODO: This is the ``Extend'' part.  Describe how you generated the new
data --- the automatic step and the human step.  Be explicit about
where the human was in the loop, because that is what makes it
\emph{semi}-automatic.  Report how much you produced against your goal:

\begin{table}[ht]
  \centering
  \begin{tabular}{lrr}
    \toprule
    & Goal & Produced \\
    \midrule
    TODO: items added          & 0 & 0 \\
    TODO: relations added      & 0 & 0 \\
    TODO: synsets touched      & 0 & 0 \\
    \bottomrule
  \end{tabular}
  \caption{What you set out to do, and what you did.}
  \label{tab:results}
\end{table}

%%% ------------------------------------------------------------------
\section{Validation}
%%% ------------------------------------------------------------------

TODO: This is the heart of the project.  Take a \emph{random} sample
(aim for at least 200 items) of what you produced and check it by hand.
Report precision.  Since you work in pairs, report agreement
(e.g.\ Cohen's $\kappa$) and say how you resolved disagreements.

The point is not the single number --- it is the \emph{categories} of
error.  Group your mistakes into types, count them, and give one real
example of each.  The categories are the most interesting thing you
will produce; they tell the next team where the method breaks.

\begin{table}[ht]
  \centering
  \begin{tabular}{lrl}
    \toprule
    Error type & Count & Example \\
    \midrule
    TODO & 0 & TODO \\
    TODO & 0 & TODO \\
    \bottomrule
  \end{tabular}
  \caption{Error types in a sample of TODO items.}
  \label{tab:errors}
\end{table}

%%% ------------------------------------------------------------------
\section{Visualisation and explanation}
%%% ------------------------------------------------------------------

TODO: One figure that makes your change --- and its errors --- legible
to someone who was not in the room.  A before/after, a coverage map, a
distribution, a sub-graph of the wordnet.  Explain what it shows and,
if your method made decisions, \emph{why} it decided as it did.

\begin{figure}[ht]
  \centering
  % \includegraphics[width=.8\linewidth]{TODO-figure}
  \fbox{\parbox[c][5cm][c]{.8\linewidth}{\centering TODO: your figure}}
  \caption{TODO: what the reader should take away from this.}
  \label{fig:main}
\end{figure}

%%% ------------------------------------------------------------------
\section{Release}
%%% ------------------------------------------------------------------

TODO: How is your contribution packaged?  It should be valid LMF that
loads with the \texttt{wn} library, with a one-line data statement and
the licence.  Say where it lives and how someone would install it.

%%% ------------------------------------------------------------------
\section{Conclusion}
%%% ------------------------------------------------------------------

TODO: What did you find?  Did you hit your number?  What did the error
analysis reveal, and what should the next team do?

%%% ------------------------------------------------------------------
\subsection*{Contribution statement}
%%% ------------------------------------------------------------------

TODO: Say who did what.  A couple of sentences is enough, for example:
``A.\ built the extraction pipeline and wrote Sections 4 and 5;
B.\ did the error analysis and the figure and wrote Sections 1--3;
both authors revised the whole paper.''

%%% ------------------------------------------------------------------
\subsection*{Use of AI tools}
%%% ------------------------------------------------------------------

TODO: If you used an LLM or other AI tool --- as part of the method or
to help write --- say which, for what, and how you checked its output.

%%% ------------------------------------------------------------------

\bibliography{report}

\end{document}
